
\section{Repository integration and migration strategy}
\label{sec:integration}

\subsection{One canonical endpoint}

The unification plan correctly identifies the ordinary
\code{Prf PRAAx} judgment as the canonical shared theorem interface.  PCFM
adopts that decision.  The native scalar calculus remains an export target;
Foundation remains an arithmetization and sentence backend; theorem packets
remain research and application layers.

A new theorem should therefore have one of the following statuses:

\begin{enumerate}
\item \textbf{Checked object theorem:}
  a committed Lean declaration with conclusion
  \(\Gamma\vdash_{\PRAAx}\varphi\).
\item \textbf{Checked conditional object theorem:}
  the same, with residual producer formulas in \(\Gamma\).
\item \textbf{Checked metatheorem:}
  a Lean theorem about translations or proof construction, explicitly
  connected to an object derivation.
\item \textbf{Executable evidence:}
  a finite checker or test that has run but is not itself an object proof.
\item \textbf{Paper proof / implementation blueprint:}
  mathematically developed but not kernel checked.
\item \textbf{Open producer:}
  a precise quantitative theorem still needed.
\end{enumerate}

These labels should appear in packet READMEs and the aggregate catalogue.

\subsection{Migration of an existing packet}

A packet should migrate in the following order.

\paragraph{Step 1: identify the endpoint.}
Write the exact \code{Fm PTm n} formula or family of formulas.  Remove semantic
phrases such as ``the encoded function has the expected meaning'' unless a
formula states that meaning.

\paragraph{Step 2: inventory finite consumers.}
List every exact computation, bounded search, finite fold, trace, certificate,
and observer.  Map each to a checked derivation or a PCFM constructor.

\paragraph{Step 3: expose producers.}
List every bound, modulus, isolation condition, finite-cover theorem, tail
estimate, or descent clock.  Give each a formula and provenance.

\paragraph{Step 4: compile one spine.}
Choose the shortest dependency path to a nontrivial endpoint and construct its
\code{Prf PRAAx} proof.  Do not first formalize every concept in the prose.

\paragraph{Step 5: discharge or retain residuals.}
Connect existing producers where available.  Publish the exact conditional
theorem for the rest.

\paragraph{Step 6: export.}
Use \code{pra_iff_native} for a minimal native proof or the bridge for a
Foundation sentence only after the ordinary theorem builds.

\subsection{Recommended first migrations}

The best initial consumers are those already closest to the checked library:

\begin{itemize}
\item exact sequence recurrences and polynomial identities;
\item finite combinatorial enumeration;
\item finite matrix calculations;
\item bounded rational approximation loops;
\item finite search theorems conditional on a supplied bound;
\item finite nerves and simplicial calculations conditional on a supplied
cover.
\end{itemize}

They exercise the module calculus without depending on unsolved analytic or
Diophantine producers.

Packets whose main claim is a global asymptotic, analytic continuation,
effective height theorem, or infinitary topological equivalence should first
extract their finite conditional core rather than attempt the whole endpoint
at once.

\subsection{CI gates}

A repository integration should add separate targets.

\begin{center}
\begin{tabularx}{\textwidth}{p{0.25\textwidth}X}
\toprule
Target & Gate\\
\midrule
\code{lake build ConciseIPRA} &
Object syntax, deduction, alternative calculi, bi-interpretations.\\
\code{lake build IPRADerivations} &
Checked arithmetic, theorem derivations, and PCFM framework.\\
\code{lake build IPRABridge} &
Both backend translations and round trips.\\
\code{python -m unittest} &
Reference algorithms and certificate-format tests; evidence only.\\
\code{latexmk} &
Research reports compile; documentation only.\\
\bottomrule
\end{tabularx}
\end{center}

The catalogue should derive checked status from CI, not from the presence of a
\code{.lean} filename.

\subsection{Metrics that measure useful progress}

Raw theorem count can reward tiny aliases and obscure the remaining boundary.
More informative metrics are:

\begin{itemize}
\item number of public \code{Prf PRAAx} endpoints;
\item number of theorem packets importing checked PCFM constructors;
\item number and type of residual obligations per endpoint;
\item size of the transitive checked dependency graph;
\item percentage of packet main claims with exact object formulas;
\item build success at the pinned commit;
\item proof-term or native-certificate size when minimality matters.
\end{itemize}

For a conditional theorem, reducing a vague ``analytic facts'' assumption to
two explicit tail inequalities is substantial progress even if the theorem is
not yet unconditional.

\subsection{Suggested directory layout}

\begin{lstlisting}[language={}]
Logic/PrimitiveRecursiveArithmetic/
  Derivations/
    Lean/IPRADerivations/
      Framework/
        Module.lean
        Boolean.lean
        ReflectQF.lean
        BoundedQuantifier.lean
        Trace.lean
        Relation.lean
        Viability.lean
        Observer.lean
        Data/
        Analysis/
        Topology/
        Algebra/
      Framework.lean
  Bridge/
    Lean/IPRABridge/
      CodeNormalization.lean
      FormulaRoundtrip.lean
      OpenSequent.lean
  docs/
    framework/
      PCFM.md
      obligation-schema.json
      migration-guide.md
\end{lstlisting}

The checked framework belongs under \code{Derivations}, not under
\code{docs/proofs}.  Research papers and uncompiled proposals remain under
documentation until they satisfy the no-placeholder rule.

\subsection{A machine-readable obligation ledger}

Each compiled theorem can emit a JSON sidecar:
\begin{lstlisting}[language={}]
{
  "theorem": "packet.namespace.endpoint",
  "object_formula": "...",
  "repository_commit": "7fa3ba48...",
  "status": "conditional-object-theorem",
  "residuals": [
    {
      "name": "tail_bound",
      "kind": "resource",
      "formula": "...",
      "source": "SpecialFunction.Tail",
      "discharged_by": null
    }
  ]
}
\end{lstlisting}

The ledger is not trusted proof data.  It supports dashboards, dependency
queries, and honest status summaries.  A CI script can verify that every
listed discharged residual corresponds to an imported theorem and that the
endpoint declaration exists.

\subsection{Immediate implementation milestones}

The smallest high-value sequence is:

\begin{enumerate}
\item commit \code{Framework/Module.lean};
\item commit Boolean truth tables and equality reflection;
\item compile the first nontrivial QF formula;
\item package \code{IPRA.ind} as fixed trace reflection;
\item port one finite recurrence packet to PCFM;
\item implement exact stable observer;
\item port one bounded-search packet with a residual bound;
\item prove code normalization for zero and composition cases;
\item finish the primitive-recursion normalization case;
\item expose the open bridge theorem.
\end{enumerate}

At each milestone there is a checked artifact that downstream work can use.
No single all-or-nothing compiler proof blocks the sequence.

\subsection{Relationship to the current CBIP and related proposals}

PCFM does not reject the mathematical content of Certified Bounded Interaction
Programs, finite-query normal forms, relational normal forms, contract DAGs,
or observation graphs.  It provides a lower implementation layer for them.

A CBIP front end may normalize to a trace and then call
\code{Framework.Trace}.  A finite-query normal form may call
\code{ReflectQF} and \code{Search}.  A relational normal form may call
\code{Viability} and \code{Observer}.  A contract DAG may elaborate every node
to \code{Module}.  This arrangement preserves the conceptual innovations of
the packets while eliminating duplication in their trusted core.
